\subsection{B-Sheaves}

Oftentimes it is more natural to define a sheaf on $X$ by defining its sections on a base of $X$, instead of on every open subset of $X$.
Recall that a base of a topology is a collection of open subsets $\B\subseteq\open(X)$ such that every open subset of $X$ can be written as
a union of sets from $\B$.

In the absence of an ambient topology, we may mean a base to mean a collection $\B$ of subsets of $X$ such that the intersection of any two
sets from $\B$ can be written as a union of elements from $\B$.
If we are given an arbitrary collection $S$ of subsets of $X$, they generate a minimal topology on $X$; $S$ is called a {\emphcolor subbase} of this
topology.
The topology generated by $\B$ consists of all finite intersections of elements from $\B$, and $\B$ is a base of this topology (in the first sense).

As $\B$ is a collection of subsets of $X$, we can view it as a poset under inclusion, and thus a subcategory of $\open(X)$.

\bdefn[title={B-Sheaves}]

    Let $X$ be a topological space with a given base $\B$.
    Define a {\emphcolor $\B$-presheaf} to be a functor $\F\colon\B^\op\to\Set$.
    This presheaf is a {\emphcolor $\B$-sheaf} if it satisfies the following two axioms:
    \benum
        \item the {\emphcolor identity axiom}: if $B=\bigcup_iB_i$ is a basic covering and $f,g\in\F(B)$ such that $f|_{B_i}=g|_{B_i}$
        for all $i$, then $f=g$.
        \item the {\emphcolor gluing axiom}: if $B=\bigcup_iB_i$ is a basic covering and $f_i\in\F(B_i)$ are given for all $i$ such that
        $f_i|_{B_k}=f_j|_{B_k}$ for all $B_k\subseteq B_i\cap B_j$, then there is an $f\in\F(B)$ such that $f|_{B_i}=f_i$ for all $i$.
    \eenum

\edefn

We can rewrite the following axioms categorically as follows.
For $B=\bigcup_iB_i$, the following diagram must be an equalizer map
$$ \F(B) \xtos\phi \prod_iB_i \xtto\alpha\beta \prod_{i,j,k}B_{i,j,k} $$
where $\set{B_{i,j,k}}_k$ ranges over all basic sets contained in $B_i\cap B_j$, and
$$ \theta = (\res_{B,B_i})_i,\qquad \alpha = (\res_{B_i,B_k}\circ\proj_i)_{i,j,k},\qquad \beta = (\res_{B_j,B_k}\circ\proj_j)_{i,j,k} . $$
Thus we can define $\B$-sheaves over any category $\C$ with enough limits.

\bdefn

    Denote the category of $\C$-valued $\B$-presheaves by $\C_\B^\pre$, and the category of $\C$-valued $\B$-sheaves by $\C_\B$.

\edefn

\bthrm

    Let $\B$ be a base for a topology on $X$, and $\F$ a sheaf of sets on this base.
    Then there is a unique sheaf of sets $\F^\ext$ on $X$ extending $\F$.

    This extends to an equivalence of categories $\ext\colon\Set_\B\xto\sim\Set_X\colon\res$, where $\res\colon\Set_X\to\Set_\B$
    is the obvious restriction functor.

\ethrm

\bproof

    For a $\B$-sheaf we can define stalks much the same was as for a usual sheaf.
    For $p\in X$ define
    $$ \F_p = \colim_{p\in B\in\B}\F(B) . $$
    Then define for $U\subseteq X$,
    $$ \F^\ext(U) = \set{(f_p\in\F_p)_{p\in U}}[\mltext{for all $p\in U$ there is a $B\in\B$ with $p\in B\subseteq U$, $s\in\F(B)$\cr
    such that $s_q=f_q$ for all $q\in B$}] . $$
    There are obvious restriction maps and it is easy to see that this is a sheaf.
    Furthermore it is easy to see that $\F^\ext(B)$ is the image of $\F(B)$ under $\iota_\F$ which is an injection as for usual sheaves.
    \qed

\eproof

\bremark

    Suppose we have an open cover $X=\bigcup_iU_i$ and for each $i$ a sheaf of sets $\F_i$ on $U_i$.
    Furthermore suppose we have isomorphisms
    $$ \phi_{ij}\colon\F_i|_{U_i\cap U_j} \xtos\sim \F_j|_{U_i\cap U_j} $$
    such that $\phi_{ii}$ is the identity and
    $$ \phi_{jk}\circ\phi_{ij} = \phi_{ik} $$
    for all $i,j,k$ on $U_i\cap U_j\cap U_k$.

    There then exists a unique sheaf $\F$ on $X$ such that $\F|_{U_i}$ is isomorphic to $\F_i$.

    This is because we can define a base $\B$ consisting of those open subsets contained in some $U_i$, i.e.\ $\B=\bigcup_i\open(U_i)$.
    Then there is a unique way to extend the $\F_i$s to a $\B$-sheaf.
    For $U\subseteq U_i$ we define $\F(U)=\F_i(U)$ (this is well-defined because if $U\subseteq U_i\cap U_j$ then $\F_i(U)\cong\F_j(U)$),
    and the restriction maps are given by the restriction maps of $\F_i$ and the $\phi_{ij}$.
    That is, for $U\subseteq U_i$ and $V\subseteq U_j$ where $U\subseteq V$ we define $\res\colon\F_j(V)\to\F_i(U)$ by
    $$ \res = \F_j(V) \xtos{\res_j} \F_j(U) \xtos{\phi_{ji}} \F_i(U) . $$
    The conditions on $\phi_{ij}$ ensure that this is a $\B$-sheaf, and we can extend it to a sheaf on $X$ via the above theorem.

\eremark

